% =============================================================================
% Section 6 — Conclusion
% =============================================================================

\section{Conclusion}
\label{sec:conclusion}

We presented NeuroAssist AI, a~privacy-preserving, offline-first assistive
communication platform that transforms residual facial gestures of patients with advanced
motor neurone disease into structured emergency alerts, delivered to caregivers over a
local network using commodity hardware.

The system's three original contributions address distinct but interrelated gaps in
current assistive technology for this population. First, by executing facial landmark
detection entirely within a~WebAssembly sandbox and serving the model binary as a local
asset, NeuroAssist AI achieves strict video-data locality without sacrificing the
accuracy or frame rate of server-side solutions. Second, the six-state deliberate gesture
state machine — which requires minimum-duration closure, confirmed inter-blink
re-opening, and a~temporal window — substantially reduces false positives relative to
simple per-frame EAR threshold comparisons, as quantified in Section~\ref{sec:evaluation}
(false-positive rate: [B]\,\% vs.\ [A]\,\% for the simple threshold; $N$=[N]
participants, [M]\,min each). Third, the hexagonal microservices architecture with
compile-time ArchUnit enforcement and the \texttt{engineFeatureRef} plug-in mechanism
provides an extensible foundation that decouples gesture detection from clinical
command management, facilitating incremental clinical expansion without architectural
regression.

End-to-end latency from gesture completion to caregiver alarm display averaged
[X]\,ms ($\sigma$=[Y]\,ms, $N$=[N] trials) on a~local network with the five-second
polling transport. This latency profile is clinically acceptable for emergency
signalling: the dominant component is the polling interval ($\leq$5\,s), not the
detection or server-processing time, and replacement with server-sent events would
reduce total latency to approximately [X]\,ms.

The hardware cost (commodity webcam, standard laptop or single-board computer, no
GPU) positions NeuroAssist AI as a~viable option for home deployment in low- and
middle-income countries, where the majority of the global MND population resides and
where specialist AAC devices are economically inaccessible. Together with the
privacy-by-design architecture, which satisfies GDPR and LGPD requirements for
biometric health data without institutional data processing infrastructure, these
properties address barriers that prevent current commercial AAC systems from reaching
patients in home care.

The primary limitation of this work is the absence of a~clinical study with MND
participants; all evaluations were performed with healthy volunteers. A~human research
ethics-approved study with patients in advanced disease stages is the essential next
step. The architecture, calibration mechanism, and measurement scripts produced in this
work are designed to support that study with minimal adaptation.
